\clearpage
\section{Herleitungen zu Abschnitt~\ref{sec:op5-coupling}}
\label{sec:anhang-b}

Dieser Anhang begründet die Gleichungen \eqref{eq:op5-drive} bis \eqref{eq:op5-drive-energy} und die anschließende Gleichung für den
Piola-Tensor. Maßgebliche Quelle ist der Companion: Die Forschungsschicht v4.10.0-r32 führt jeden Schritt mit Beweis
\cite{companion,repo}; dieser Anhang gibt sie auf Deutsch wieder. Jeder Rechenschritt ist exakt mit Wolfram Language nachgeprüft; Skript und Originalausgabe liegen im Quellrepository
unter \texttt{paper/\allowbreak preprint-de/\allowbreak audit/\allowbreak anhang-b/}. Zuerst wird offengelegt, welche Gleichungen keine Herleitung haben, sondern Annahmen
sind.

\subsection{Was angenommen und was hergeleitet ist}
\label{sec:anhang-b-annahmen}

Zwei Eingaben des Abschnitts sind gesetzt.
\begin{itemize}
 \item \textbf{Der Antriebsoperator \eqref{eq:op5-drive}.} Seine Form folgt aus dem reduzierten OP5-Vergleichsmodell mit einem
 gerichteten massiven Vermittler und dem Einstein-Kanal; der Einstein-Beitrag \(\tfrac32 g_0^2 d\) ist im Haupttext an der
 transversalen Verbesserung \(S^{\mu\nu}\) nachvollzogen. Die Koeffizienten \(Z\), \(r_m\), \(c_h\), \(g_0\), \(g_m\), \(\Omega_0\) und \(m\) sind
 dagegen gewählt. Die im Text genannte Wahl ist eine in den Notizen geprüfte, keine abgeleitete.
 \item \textbf{Die Gleichsetzung \eqref{eq:op5-match}.} Sie setzt die radiale Higgsquelle \(\delta\mathcal J_h\) der Materie mit dem
 Antrieb \(\delta J_{\rm OP5}\) des Vergleichsfelds gleich, in der Normierung, in der beide dieselbe Einheit tragen. Begründet wäre sie
 durch eine Wirkung, in der der Higgsradius \(h_{\rm W}\) mit dem Koeffizienten eins an das Vergleichsfeld \(h\) koppelt, also durch einen
 Term der Form \(\int h\,\mathcal J_h\) in derselben Wirkung, aus der auch \eqref{eq:op5-drive} folgt. Eine solche Wirkung liegt nicht vor.
\end{itemize}
Die Gleichsetzung trägt nur die gerichtete Hälfte der Rechnung. Ohne sie gelten weiterhin die mechanische Konstruktion, die Erhaltung
des präparierten Tensors \eqref{eq:op5-tensor} und die Kompaktheit des gewöhnlichen Außenbeitrags \eqref{eq:op5-compact}; der Tensor
\(\delta T\) hängt nur von der Verschiebung ab. Ohne sie entfällt dagegen die Aussage, dass dieselbe Präparation den gerichteten Kontrast
\eqref{eq:op5-contrast} antreibt. Mit einer anderen Normierung \(\delta J_{\rm OP5}=c\,\delta\mathcal J_h\) skaliert der Kontrast mit \(c\);
seine Form und sein Vorzeichen bleiben. Ob \(c\neq0\) ist, entscheidet erst eine gemeinsame Wirkung.

Eine Prüfung am Higgs-Sektor der Wirkung W2 verschiebt die Frage. Ist das Vergleichsfeld der Radialmodus \(\eta\) des Higgs-Dubletts, so gilt
\(2(\partial_t^2-\Delta)\eta+4\mu^2\eta=\mathcal J_h\), und die Gleichsetzung folgt für den skalaren Kanal exakt. Diese Lesart erzwingt aber
\(c_h=1\) statt der gesetzten \(1/2\), und bei minimaler Kopplung gibt es weder eine Einstein-Mischung (\(g_0=0\)) noch einen linear
gekoppelten gerichteten Vermittler (\(g_m=0\)). Der gerichtete Kontrast beruht damit auf einer anderen Annahme: einem unterluminalen
massiven Vermittler mit linearer Kopplung an den Higgsradius, den W2 nicht enthält. Quelle ist der Companion, Runde v4.10.0-r33
\cite{companion,repo}. Eine Suche im Companion (Runde v4.10.0-r34) findet kein Vakuumfeld mit diesen Eigenschaften. Der gegenläufige
Paar-Konnektor aus einem Materiezustand und seinem Tachyonbild ist massiv, unterluminal und gerichtet, hat aber weder Wirkung noch eine
belegte Kopplung an den Higgsradius. Die optische Mode des zweikomponentigen Materials erfüllt alle Anforderungen; mit ihr wäre der
gerichtete Kontrast eine Eigenschaft der Präparation. Runde v4.10.0-r35 rechnet diesen Weg durch: Bei einachsiger Präparation entsteht ein
gerichteter Term mit der optischen Lücke als Masse und dem Gewicht \(\mu(Y_1/\rho_1-Y_2/\rho_2)^2\). Er koppelt nur, wenn die beiden Komponenten
verschieden große Higgs-Anteile ihrer Ruhemasse haben. Das führende Glied des Kontrasts braucht \(r_m\) nicht; der Zähler ist aber \(a_\perp\),
eine Potenz weniger als im Vergleichsmodell.

\subsection{Die Quelle der elastischen Verschiebung}
\label{sec:anhang-b-quelle}

Für einen Gradienten \(\mathbf b=\nabla S\) gilt mit dem isotropen Cauchy-Tensor
\[
 \sigma(\nabla S)=\lambda(\Delta S)\,I+2\mu\,\nabla\nabla S,\qquad
 \mathcal L\nabla S=\nabla\cdot\sigma(\nabla S)=(\lambda+2\mu)\nabla\Delta S=K_L\nabla\Delta S .
\]
Daraus folgt \(\nabla\cdot\mathcal L\mathbf b=K_L\Delta^2S\). Mit \eqref{eq:op5-tensor} ist \(\delta T^{00}=-\beta_\rho K_L\Delta^2S\), und mit
\eqref{eq:op5-density}, \(\delta\epsilon_{\rm mat}|_v=-v\,\delta\mathcal J_h\), ergibt sich \eqref{eq:op5-elas-source}:
\(\delta\mathcal J_h=-\delta T^{00}/v=(\beta_\rho K_L/v)\Delta^2S\).

\subsection{Die Wahl des Profils, Gleichung \eqref{eq:op5-profile}}
\label{sec:anhang-b-profil}

Setzt man \(S=-\dfrac{v}{\beta_\rho K_L}D_0\mathcal P_{\rm dyn}V\) in \eqref{eq:op5-elas-source} ein, so ist
\[
 \delta\mathcal J_h=-\Delta^2D_0\mathcal P_{\rm dyn}V=\mathcal P_{\rm dyn}\bigl(-D_0\Delta^2V\bigr)=\mathcal P_{\rm dyn}\,\delta H ,
\]
weil Operatoren mit konstanten Koeffizienten kommutieren. Das ist genau die Antwort, die das Vergleichsmodell auf \(\delta H\) verlangt,
sodass \eqref{eq:op5-match} auf diesem Profil erfüllt ist. Zugleich ist \(\mathbf b=\nabla S=-D_0\mathbf b_*\) mit
\(\mathbf b_*=(v/(\beta_\rho K_L))\nabla\mathcal P_{\rm dyn}V\). Da \(\delta T\) in \eqref{eq:op5-tensor} linear in \(\mathbf b\) ist und nur
Ableitungen mit konstanten Koeffizienten enthält, folgt \(\delta T[\mathbf b]=-D_0\,T_*\) mit \(T_*=\delta T[\mathbf b_*]\). \(T_*\) ist als Tensor
derselben Form wie \eqref{eq:op5-tensor} für sich erhalten.

\subsection{Der kompakte gewöhnliche Außenbeitrag, Gleichung \eqref{eq:op5-compact}}
\label{sec:anhang-b-kompakt}

In harmonischer Eichung lautet die linearisierte Einstein-Gleichung \(D_0\bar\gamma_{\mu\nu}=16\pi G\,\delta T_{\mu\nu}\) mit
\(D_0=\partial_t^2-\Delta\). Setzt man \(\delta T=-D_0T_*\) ein, ist
\[
 D_0\bigl(\bar\gamma_{\mu\nu}+16\pi G\,T_{*,\mu\nu}\bigr)=0 .
\]
Vor der Schaltung verschwinden \(V\), \(T_*\) und \(\bar\gamma\); eine einlaufende freie Welle ist ausgeschlossen. Die retardierte Lösung der
homogenen Wellengleichung mit verschwindenden Anfangsdaten ist null, also \(\bar\gamma_{\mu\nu}=-16\pi G\,T_{*,\mu\nu}\). Weil \(V\) räumlich
kompakt ist, ist es auch \(T_*\), und damit verschwinden \(\bar\gamma\) und seine lineare Krümmung außerhalb dieses Trägers, auch
während der Schaltung. Die Aussage benutzt nur die Linearität, die konstanten Koeffizienten und die Anfangsdaten; die Gleichsetzung
\eqref{eq:op5-match} geht nicht ein.

\subsection{Das Achsenprofil, Gleichungen \eqref{eq:op5-axis-profile} und \eqref{eq:op5-normalized-axis}}
\label{sec:anhang-b-achse}

Mit \(\xi=r_\perp^2/R^2\) und dem zweidimensionalen radialen Laplace-Operator \(\Delta_\perp=\partial_r^2+r^{-1}\partial_r\) gilt für \(r_\perp^{2j}\)
\[
 \Delta_\perp\,\xi^j=\frac{(2j)^2}{R^2}\,\xi^{j-1}=\frac{4j^2}{R^2}\,\xi^{j-1}.
\]
Dreimal angewandt ergibt das für die beiden Terme des Profils die Faktoren \(4^3\cdot25\cdot16\cdot9=230400\) und
\(4^3\cdot36\cdot25\cdot16=921600=4\cdot230400\), also
\[
 \Delta_\perp^3\bigl(\xi^5-\kappa\xi^6\bigr)=\frac{230400}{R^6}\bigl(\xi^2-4\kappa\xi^3\bigr).
\]
Der Vorfaktor \(-\mathcal A R^6/230400\) in \eqref{eq:op5-axis-profile} ist so gewählt, dass \(\Delta_\perp^3V=-\mathcal A\,p(z)(\xi^2-4\kappa\xi^3)\)
im flachen Kern, wo \(\chi=1\) ist. Auf dem statischen Plateau ist \(D_0=-\Delta\) und damit \(\delta H=-D_0\Delta^2V=\Delta^3V\). Schreibt man
\(\Delta^3=(\Delta_\perp+\partial_z^2)^3\), enthält jeder Term mit einer Längsableitung höchstens zwei radiale Laplace-Operatoren; sein
niedrigster Grad in \(\xi\) ist \(\xi^3\). Daher ist \(\delta H=-\mathcal A p(z)\,\xi^2+O(\xi^3)\). Mit \(\Delta_\perp\xi^2=16\xi/R^2\) folgt die zweite
Form in \eqref{eq:op5-normalized-axis}. Die Wahl von \(\kappa\) als Quotient der Momente sorgt dafür, dass das Querintegral
\(\int V\,d^2r_\perp=\pi R^2\int_0^\infty V\,d\xi\) je Längsschnitt verschwindet.

\subsection{Der Gezeitenkontrast, Gleichung \eqref{eq:op5-contrast}}
\label{sec:anhang-b-kontrast}

Im massiven Vergleich ist \(f=\Delta_\perp H/m^2\), also nahe der Achse
\[
 \delta f=\frac{\Delta_\perp\delta H}{m^2}=-\frac{16\mathcal A p(z)}{m^2R^2}\,\xi+O(\xi^2).
\]
Für eine transversale Einheitsrichtung \(e\) ist \(\partial_e^2\xi=2/R^2\), und Terme ab \(\xi^2\) haben auf der Achse keine zweite Ableitung.
Damit ist \(\partial_e^2(\delta f)(0,z)=-32\mathcal A p(z)/(m^2R^4)\). Die verwendete Auslese ist der geometrisch-optische Gezeitenterm eines
Strahls entlang der Achse, \(\Delta\mathcal I=-\tfrac12\int\partial_e^2(\delta f)\,dz\); er misst die integrierte transversale Krümmung des
Vergleichsfelds für eine feste Polarisation. Einsetzen gibt \(\Delta\mathcal I=16\mathcal A\int p\,dz/(m^2R^4)\). Die Normierung
\(-\tfrac12\) ist eine Konvention der Auslese; das Vorzeichen und der Faktor hängen an ihr, das Nicht-Verschwinden nicht.

\subsection{Der Bündelfaktor, Gleichung \eqref{eq:op5-beam}}
\label{sec:anhang-b-buendel}

Integriert man über \(z\), fallen alle Beiträge mit einer Längsableitung als Randterme weg, weil \(p\) kompakt ist; es bleibt
\(\int\delta f\,dz=-\mathcal A P\,\Delta_\perp(\xi^2-4\kappa\xi^3)/m^2\) mit \(P=\int p\,dz\). Für eine radiale Funktion \(F(r)\) ist das Mittel
von \(\partial_e^2F\) über den Azimut eines Strahls im Abstand \(r\) gleich \(\tfrac12\Delta_\perp F\), weil
\(\partial_x^2F=F''\cos^2\phi+(F'/r)\sin^2\phi\). Mit \(\Delta_\perp\xi=4/R^2\), \(\Delta_\perp\xi^2=16\xi/R^2\) und \(\Delta_\perp\xi^3=36\xi^2/R^2\) folgt
für das Verhältnis zum Achsenwert
\[
 \frac{\langle\partial_e^2\delta f\rangle_\phi(\xi)}{\partial_e^2\delta f(0)}=1-36\kappa\,\xi .
\]
Über ein achsenparalleles Bündel mit rotationssymmetrischem Gewicht gemittelt, entsteht \eqref{eq:op5-beam} mit \(\langle\xi\rangle\).

\subsection{Die planare Kontrolle, Gleichung \eqref{eq:op5-planar-check}}
\label{sec:anhang-b-planar}

Der Kern \(K_m(z)=\tfrac m2e^{-m|z|}\) hat das Integral eins; also ist \(\int_{\mathbb R}f\,dz=\int\delta\rho\,dz\). Liegt der Träger von
\(\delta\rho\) in \([0,L]\), ist für \(z<0\)
\(f(z)=\int\tfrac m2e^{-m(z'-z)}\delta\rho(z')\,dz'=e^{mz}f(0)\), und entsprechend \(f(z)=e^{-m(z-L)}f(L)\) für \(z>L\). Die beiden
Ausläufer haben die Integrale \(f(0)/m\) und \(f(L)/m\). Zieht man sie vom Gesamtintegral ab, steht \eqref{eq:op5-planar-check} da.

\subsection{Die Energiebilanz beim Schalten, Gleichung \eqref{eq:op5-drive-energy}}
\label{sec:anhang-b-energie}

Die mechanische Energiedichte ist
\[
 e_{\rm mech}=\tfrac12\rho_1|\dot{\mathbf u}|^2+\tfrac12\rho_2|\dot{\mathbf w}|^2+W(\varepsilon(\mathbf u))+c_\sigma W(\varepsilon(\mathbf w))
 +\tfrac k2|\mathbf R_{\rm F}|^2 ,
\]
mit der elastischen Energie \(W\) zu \(\mathsf C\). Ihre Zeitableitung enthält \(\partial_tW(\varepsilon(\mathbf u))=\sigma_1:\nabla\dot{\mathbf u}
=\nabla\cdot(\sigma_1\dot{\mathbf u})-\dot{\mathbf u}\cdot\mathcal L\mathbf u\), entsprechend für \(\mathbf w\). Auf den Bewegungsgleichungen
\(\rho_1\ddot{\mathbf u}=\mathcal L\mathbf u-k\mathbf R_{\rm F}\) und \(\rho_2\ddot{\mathbf w}=c_\sigma\mathcal L\mathbf w+k\mathbf R_{\rm F}\) heben sich
die elastischen Volumenterme gegen die kinetischen, und die Federterme ergeben
\(-k\mathbf R_{\rm F}\cdot\dot{\mathbf u}+k\mathbf R_{\rm F}\cdot\dot{\mathbf w}+k\mathbf R_{\rm F}\cdot(\dot{\mathbf u}-\dot{\mathbf w}-\dot{\mathbf d})=-k\,\mathbf R_{\rm F}\cdot\dot{\mathbf d}\).
Übrig bleiben \(\nabla\cdot\mathbf F\) und \(P_{\rm Schalt}\). Die Wolfram-Prüfung rechnet den eindimensionalen Fall mit derselben Struktur.

\subsection{Der Piola-Tensor in zweiter Ordnung}
\label{sec:anhang-b-piola}

Für das Saint-Venant--Kirchhoff-Gesetz ist \(P=(I+\mathsf H)\,\mathsf C\,E\) mit \(E=\tfrac12(\mathsf H+\mathsf H^{\mathsf T}+\mathsf H^{\mathsf T}\mathsf H)\).
Der quadratische Anteil in \(\mathsf H\) ist \(\mathsf C(\tfrac12\mathsf H^{\mathsf T}\mathsf H)+\mathsf H\,\mathsf C\,\varepsilon\). Für symmetrisches
\(\mathsf H\) ist \(\mathsf H^{\mathsf T}\mathsf H=\mathsf H^2\) und
\[
 \mathsf C(\tfrac12\mathsf H^2)+\mathsf H\,\mathsf C\,\mathsf H
 =\tfrac\lambda2\operatorname{tr}(\mathsf H^2)I+\mu\mathsf H^2+\lambda(\operatorname{tr}\mathsf H)\mathsf H+2\mu\mathsf H^2=N(\mathsf H).
\]
Die zweite Komponente mit Steifigkeit \(c_\sigma\mathsf C\) und \(\mathbf w_1=-\mathbf u_1/c_\sigma\) trägt \(c_\sigma N(-\mathsf H/c_\sigma)=N(\mathsf H)/c_\sigma\),
weil \(N\) quadratisch ist. Die linearen Anteile zweiter Ordnung addieren sich zu \(\mathsf C\varepsilon(\mathbf u_2)+c_\sigma\mathsf C\varepsilon(\mathbf w_2)
=\mathsf C\varepsilon(S_2)\). Zusammen ergibt das die Gleichung für \((P_1+P_2)^{(2)}\).

\subsection{Was die Prüfung abdeckt}
\label{sec:anhang-b-pruefung}

Exakt geprüft sind: die Identität \(\mathcal L\nabla S=K_L\nabla\Delta S\) und \(\nabla\cdot\mathcal L\nabla S=K_L\Delta^2S\); die Faktoren \(4j^2\) und
230400; der Achsenwert \(16\mathcal A P/(m^2R^4)\); der Bündelfaktor \(1-36\kappa\xi\); die planare Identität an einem Beispielprofil
(sie ist linear in \(\delta\rho\)) samt Ausläufern und Gesamtintegral; die Energiebilanz im eindimensionalen Fall; der quadratische
Piola-Anteil \(N(\mathsf H)\) und der Faktor \(1/c_\sigma\). Nicht geprüft und nicht prüfbar ohne eine gemeinsame Wirkung sind die Koeffizienten von
\eqref{eq:op5-drive} und die Gleichsetzung \eqref{eq:op5-match}.
